\begin{abstract}
Neural receivers outperform conventional 5G~NR processing chains, but their
compute and memory demands hinder real-time deployment. For a
standard-compliant \gls{MU}-\gls{MIMO} neural receiver, the 4-bit number
\emph{format}, not merely the bit width, determines whether compression
preserves the gain over classical receivers. We apply weight and activation
\gls{QAT} and, separately, $50\%$ magnitude pruning, comparing
INT8/INT4 and FP8~(E4M3)/FP4~(E2M1) weights with INT8 post-\gls{ReLU}
activations. Trained on 3GPP UMi channels and evaluated on TDL-B and TDL-C,
8-bit weight--activation models remain within $0.05$\,dB of \texttt{FP32} at
$10\%$ and $1\%$ \gls{BLER}. At 4~bits, uniform INT4 loses
$3.3$--$3.7$\,dB and falls below LS--\gls{LMMSE}, whereas FP4 more than
halves this loss ($1.3$--$1.4$\,dB) and still outperforms it by about
$0.5$\,dB, even after pruning. FP4's denser near-zero grid matches the
trained weight distribution, and FP4 avoids the residual-path over-pruning
seen with INT4. An analytic cost model projects $66\times$ fewer
bit-operations and $8.8\times$ less weight storage for pruned 4-bit-weight
inference.
\end{abstract}

%
\begin{keywords}
Neural receivers, MU-MIMO, quantization-aware training, FP4, pruning
\end{keywords}

%
\begin{comment}
    for IEEE submission portal :
Neural receivers outperform conventional 5G NR processing chains, but their compute and memory demands hinder real-time deployment. For a standard-compliant multi-user MIMO neural receiver, the 4-bit number format, not merely the bit width, determines whether compression preserves the gain over classical receivers. We apply weight and activation quantization-aware training (QAT) and, separately, 50% magnitude pruning, comparing INT8/INT4 and FP8 (E4M3)/FP4 (E2M1) weights with INT8 post-ReLU activations. Trained on 3GPP UMi channels and evaluated on TDL-B and TDL-C, 8-bit weight-activation models remain within 0.05 dB of FP32 at 10% and 1% block error rate (BLER). At 4 bits, uniform INT4 loses 3.3-3.7 dB and falls below LS-LMMSE, whereas FP4 more than halves this loss (1.3-1.4 dB) and still outperforms it by about 0.5 dB, even after pruning. FP4's denser near-zero grid matches the trained weight distribution, and FP4 avoids the residual-path over-pruning seen with INT4. An analytic cost model projects 66x fewer bit-operations and 8.8x less weight storage for pruned 4-bit-weight inference.
\end{comment}
%
